\chapter{Tổng quan hệ thống}

\section{Bài toán và động cơ}

LexiLingo là nền tảng học tiếng Anh theo mô hình \en{gamified microlearning} (học chia
nhỏ có yếu tố trò chơi) tương tự Duolingo, nhưng bổ sung một lớp \textbf{cá nhân hoá dựa
trên đồ thị tri thức người học}. Bài toán trung tâm mà hệ thống giải quyết không phải là
"trình bày nội dung tiếng Anh", mà là:

\begin{quote}
\itshape Làm sao để mọi hoạt động luyện tập của người học — làm bài, chơi game, ôn từ,
nói chuyện với AI — đều trở thành \textbf{bằng chứng đo lường được} về năng lực, và
được dùng ngay để quyết định nội dung kế tiếp cũng như nội dung phản hồi của trợ lý AI?
\end{quote}

Hầu hết ứng dụng học ngôn ngữ dừng ở mức cộng điểm kinh nghiệm và đếm bài đã học.
LexiLingo đi xa hơn ở ba điểm:

\begin{enumerate}[leftmargin=1.4em]
  \item \textbf{Hai động cơ mô hình hoá người học chạy song song.} Một động cơ chấm
  điểm kỹ năng theo khung CEFR (\en{Common European Framework of Reference} — Khung
  tham chiếu ngôn ngữ chung châu Âu), một động cơ lập lịch ôn tập theo trạng thái trí
  nhớ của từng khái niệm. Mọi hoạt động có chấm điểm đều phải chảy vào cả hai.
  \item \textbf{Trợ lý AI được "nối đất" (\en{grounded}) vào trạng thái người học} thay
  vì chỉ vào tài liệu tĩnh. Đây là kiến trúc TRACE-CAG được trình bày ở Chương~8.
  \item \textbf{Nội dung khoá học sinh bởi tác tử AI} (\en{content agent}) có kiểm duyệt,
  thay vì nhập tay toàn bộ.
\end{enumerate}

\section{Phạm vi hệ thống}

Hệ thống là một \en{monorepo} (kho mã nguồn hợp nhất — nhiều dịch vụ trong cùng một kho
Git) gồm sáu thành phần triển khai được và một số thành phần hỗ trợ:

\begin{longtable}{L{3.2cm} L{4.2cm} L{6.2cm}}
\toprule
\textbf{Thành phần} & \textbf{Công nghệ} & \textbf{Trách nhiệm}\\
\midrule
\endhead
\code{flutter-app/} & Flutter 3.x, Dart \(\ge\)3.8.1, Provider, GetIt &
Ứng dụng người học đa nền tảng (Android, iOS, Web). Phiên bản hiện tại 1.7.0+17, 26 mô-đun tính năng.\\
\code{backend-service/} & FastAPI, SQLAlchemy 2.0 \en{async}, PostgreSQL, Redis, Celery &
Toàn bộ nghiệp vụ: xác thực, khoá học, tiến độ, từ vựng, trò chơi, gamification, đánh giá năng lực, quản trị.\\
\code{ai-service/} & FastAPI, LangGraph, KuzuDB, MongoDB, Redis &
Pipeline TRACE-CAG, hội thoại, giọng nói (STT/TTS/phát âm), dịch, sinh nội dung, tác tử xếp hạng và thông báo.\\
\code{admin-service/} & React 19, TypeScript, Vite, Ant Design &
Bảng điều khiển quản trị: quản lý nội dung, người dùng, mô hình AI, giám sát.\\
\code{gateway/} & Nginx 1.27 (kèm cấu hình Kong, Cloudflare WAF) &
Cổng API duy nhất: định tuyến, giới hạn tần suất, chống tấn công biên.\\
\code{mcp-server/} & Model Context Protocol &
Công cụ cho lập trình viên và tác tử AI trong IDE (quản lý khoá i18n, gọi mô hình cục bộ).\\
\bottomrule
\caption{Sáu thành phần triển khai của LexiLingo}
\end{longtable}

Ngoài ra: \code{contracts/} (hợp đồng dữ liệu dùng chung giữa các dịch vụ),
\code{monitoring/} (Prometheus + Grafana), \code{deploy/} và \code{scripts/}
(triển khai, sao lưu, kiểm thử khói), \code{system-testing/} (kiểm thử liên dịch vụ).

\section{Giải thích thuật ngữ tiếng Anh cốt lõi}

Báo cáo dùng nhiều thuật ngữ tiếng Anh vì đó là tên gọi trong mã nguồn. Bảng dưới giải
thích các thuật ngữ xuất hiện xuyên suốt; phụ lục A liệt kê đầy đủ hơn.

\begin{longtable}{L{3.6cm} L{10.2cm}}
\toprule
\textbf{Thuật ngữ} & \textbf{Giải thích}\\
\midrule
\endhead
\en{CEFR} & Khung tham chiếu ngôn ngữ chung châu Âu, sáu bậc A1--C2. LexiLingo dùng làm thang đo năng lực chính.\\
\en{LLM} (\en{Large Language Model}) & Mô hình ngôn ngữ lớn — mạng nơ-ron sinh văn bản, ví dụ Qwen, Gemini, Llama.\\
\en{RAG} (\en{Retrieval-Augmented Generation}) & Sinh văn bản có truy xuất: tìm tài liệu liên quan rồi đưa vào ngữ cảnh cho LLM.\\
\en{CAG} (\en{Cache-Augmented Generation}) & Sinh văn bản có bộ nhớ đệm: tái sử dụng ngữ cảnh/kết quả đã tính thay vì truy xuất lại từ đầu.\\
\en{Knowledge Graph} (KG) & Đồ thị tri thức — biểu diễn khái niệm là đỉnh, quan hệ (tiên quyết, liên quan, dễ nhầm) là cạnh.\\
\en{Pipeline} & Chuỗi bước xử lý nối tiếp, đầu ra bước trước là đầu vào bước sau.\\
\en{StateGraph} & Đồ thị trạng thái của LangGraph: mỗi đỉnh là một hàm biến đổi trạng thái, cạnh quyết định bước kế tiếp.\\
\en{FSRS} (\en{Free Spaced Repetition Scheduler}) & Bộ lập lịch lặp lại ngắt quãng thế hệ mới, mô hình hoá trí nhớ bằng \en{stability} và \en{difficulty}.\\
\en{Spaced repetition} & Lặp lại ngắt quãng — ôn lại ngay trước thời điểm dự đoán sẽ quên.\\
\en{Retrievability} & Xác suất nhớ lại tại một thời điểm, ký hiệu \(R(t)\).\\
\en{Stability} & Độ bền trí nhớ, tính bằng ngày; càng lớn thì quên càng chậm.\\
\en{Mastery} & Mức độ nắm vững một khái niệm, giá trị trong \([0,1]\).\\
\en{EMA} (\en{Exponential Moving Average}) & Trung bình trượt luỹ thừa — cập nhật điểm mới bằng cách pha trộn có trọng số với điểm cũ.\\
\en{Idempotent} & Bất biến khi lặp — gọi nhiều lần cho cùng kết quả như gọi một lần, quan trọng khi mạng gửi lại yêu cầu.\\
\en{Outbox} & Mẫu thiết kế "hộp thư đi": ghi sự kiện vào cùng giao dịch cơ sở dữ liệu rồi mới phát đi, tránh mất sự kiện.\\
\en{Spool} & Vùng đệm ghi trước — hàng đợi bền vững lưu quan sát trước khi chuyển tiếp.\\
\en{Middleware} & Lớp trung gian xử lý mọi yêu cầu HTTP trước/sau khi vào ứng dụng.\\
\en{RBAC} (\en{Role-Based Access Control}) & Phân quyền theo vai trò.\\
\en{ETL} (\en{Extract–Transform–Load}) & Rút trích – biến đổi – nạp dữ liệu từ nguồn ngoài vào hệ thống.\\
\en{STT / TTS} & \en{Speech-to-Text} (nhận dạng giọng nói) / \en{Text-to-Speech} (tổng hợp giọng nói).\\
\en{VAD} (\en{Voice Activity Detection}) & Phát hiện có tiếng nói, dùng để biết khi nào người học bắt đầu/dừng nói.\\
\en{SSE} (\en{Server-Sent Events}) & Kênh một chiều máy chủ đẩy dữ liệu về trình duyệt, dùng cho phản hồi AI theo dòng.\\
\en{Latency} & Độ trễ; \en{p95/p99} là phân vị 95\%/99\% của độ trễ.\\
\en{Throughput} & Thông lượng — số yêu cầu xử lý được mỗi đơn vị thời gian.\\
\bottomrule
\caption{Thuật ngữ tiếng Anh cốt lõi}
\end{longtable}

\section{Con số tổng quan}

Các con số dưới đây được đếm trực tiếp từ mã nguồn tại thời điểm viết báo cáo
(nhánh \code{dev}, tháng 8/2026):

\begin{longtable}{L{7.2cm} L{6.4cm}}
\toprule
\textbf{Hạng mục} & \textbf{Số lượng}\\
\midrule
\endhead
Mô-đun định tuyến backend (\code{app/routes/}) & 41 tệp\\
Mô hình ORM backend (\code{app/models/}) & 25 tệp, khoảng 70 bảng\\
Dịch vụ nghiệp vụ backend (\code{app/services/}) & 45 tệp\\
Tác vụ nền Celery (\code{app/tasks/}) & 9 mô-đun, 12 lịch định kỳ\\
Mô-đun định tuyến AI service & 17 tệp\\
Dịch vụ AI (\code{api/services/}) & 45 tệp + 5 gói con\\
Mã nguồn TRACE-CAG (\code{trace\_cag/}) & 8.078 dòng trên 19 tệp\\
Mô-đun tính năng Flutter (\code{lib/features/}) & 26 mô-đun\\
Phụ thuộc Flutter & 63 gói\\
Trang quản trị React & hơn 25 trang\\
Bộ kiểm thử & 102 tệp (backend), 61 (AI), 120 (Flutter)\\
\bottomrule
\caption{Quy mô mã nguồn}
\end{longtable}

\section{Cấu trúc báo cáo}

Chương~2 phân tích hệ thống ở góc nhìn nghiệp vụ: tác nhân, use case, quy trình và
quy tắc ràng buộc. Chương~3 trình bày kiến trúc tổng thể và luồng yêu cầu. Chương~4 là
thiết kế hệ thống: use case hệ thống, ERD, sơ đồ lớp, sơ đồ tuần tự và sơ đồ triển khai.
Chương~5 đi sâu vào backend-service, Chương~6 mô tả toàn bộ tầng dữ liệu (PostgreSQL,
Redis, MongoDB, KuzuDB). Chương~7 giới thiệu AI service ở mức tổ chức; Chương~8 phân
tích sâu thuật toán TRACE-CAG — từng đỉnh xử lý, thuật toán mở rộng đồ thị, điều khiển
thích nghi, IRCoT. Chương~9 là chương thuật toán nền tảng: công thức đầy đủ của FSRS
delta-rule, chấm điểm CEFR, bộ nhớ đệm SCAR-L1, hoà trộn truy xuất và chống cày điểm.
Chương~10 mô tả các pipeline dữ liệu. Chương~11 tổng hợp API. Chương~12 và~13 trình
bày phía người dùng (Flutter, admin web, gateway). Chương~14 nói về kiểm thử, triển
khai và vận hành. Chương~15 là phụ lục thuật ngữ và bảng dữ liệu.
